% !PRINTED 485
\weberpage{485}
%<<CONT
For what still remains is the proof that one can dispose of $\eta_0, \eta_0', \eta_0''$ (§.\,130) in such a way that the functions $X_1, X_2, X_3$ do not vanish identically. For this we must consider the mode of formation of the quantities $X$ somewhat more closely.

If, in accordance with (8) (§.\,130), we set in (9):
\[
\frac{\partial^3 f}{\partial \xi_1^3} = 6\,\xi_3, \quad \frac{\partial^3 f}{\partial \xi_1^2 \xi_2} = 0, \quad \frac{\partial^3 f}{\partial \xi_1^2 \partial \xi_3} = 6\,\xi_1, \quad \frac{\partial^3 f}{\partial \xi_1 \partial \xi_2 \partial \xi_3} = 0, \dots, % [sic] source omits the second \partial before \xi_2
\]
then it results that

\weq{1}{\begin{aligned}
X_1 &= A_{2,2,2} + 3\,A_{1,1,3}, \quad X_2 = A_{3,3,3} + 3\,A_{2,2,1}, \\
X_3 &= A_{1,1,1} + 3\,A_{3,3,2}.
\end{aligned}}

Now, however, by (4) and (6) (§.\,130) we further have:

\weq{2}{\Psi =
\begin{vmatrix}
p_1 x_1 + q_3 x_2 + q_2 x_3, & q_3 x_1 + p_2 x_2 + q_1 x_3, & q_2 x_1 + q_1 x_2 + p_3 x_3 \\
p_1' x_1 + q_3' x_2 + q_2' x_3, & q_3' x_1 + p_2' x_2 + q_1' x_3, & q_2' x_1 + q_1' x_2 + p_3' x_3 \\
p_1'' x_1 + q_3'' x_2 + q_2'' x_3, & q_3'' x_1 + p_2'' x_2 + q_1'' x_3, & q_2'' x_1 + q_1'' x_2 + p_3'' x_3
\end{vmatrix}.}

This can easily be ordered according to powers and products of the $x$, and if we accordingly use the notation
\[
(p_1, q_3, q_2) =
\begin{vmatrix}
p_1, & q_3, & q_2 \\
p_1', & q_3', & q_2' \\
p_1'', & q_3'', & q_2''
\end{vmatrix} \cdots,
\]
one obtains
\[
\begin{aligned}
A_{1,1,1} &= (p_1, q_3, q_2), & 3\,A_{3,3,2} &= (q_3, q_1, p_3) + (q_2, p_2, p_3) \\
A_{2,2,2} &= (q_3, p_2, q_1), & 3\,A_{1,1,3} &= (p_1, q_1, q_2) + (p_1, q_3, p_3) \\
A_{3,3,3} &= (q_2, q_1, p_3), & 3\,A_{2,2,1} &= (q_3, p_2, q_2) + (p_1, p_2, q_1),
\end{aligned}
\]
and from this, by (1),

\weq{3}{\begin{aligned}
X_1 &= (q_3, p_2, q_1) + (p_1, q_1, q_2) + (p_1, q_3, p_3) \\
X_2 &= (q_2, q_1, p_3) + (q_3, p_2, q_2) + (p_1, p_2, q_1) \\
X_3 &= (p_1, q_3, q_2) + (q_3, q_1, p_3) + (q_2, p_2, p_3).
\end{aligned}}

These functions $X_1, X_2, X_3$ are, as follows from what was remarked above [§.\,130, (5)], trilinear forms of the three sequences of variables $\eta_r, \eta_r', \eta_r''$, whose coefficients can be expressed rationally in terms of the seventh root of unity $\varepsilon$. But the coefficients of these forms are certainly not all equal to zero. For by (3) the quantities $p, q$ can be so chosen that $X_1, X_2, X_3$ do not become equal to zero, and the 21 quantities $\eta_r, \eta_r', \eta_r''$ can be determined from §.\,130, (5) in such a way that the 18 quantities $p, q$ (and moreover the three sums $\Sigma \eta_r$) obtain arbitrarily prescribed values
%CONT>>
